% !TEX program = xelatex
\documentclass[10pt,letterpaper,twocolumn]{article}

\usepackage[top=0.85in, bottom=0.85in, left=0.75in, right=0.75in]{geometry}
\usepackage{fontspec}
\usepackage{unicode-math}
% Latin Modern has NO Cyrillic glyphs. STIX Two Text is a Latin-Modern-
% adjacent serif with full Cyrillic that ships on macOS AND in TeX Live
% (stix2-otf), so this preamble renders identically on macOS xelatex and
% the oven's Linux xelatex.
\setmainfont{STIX Two Text}
\setsansfont{STIX Two Text}
% YWFT Processing is Latin-only (no Cyrillic); \acbold is kept defined for
% the Latin title "The Pals Mark" only.
\newfontfamily\acbold{ywft-processing-bold}[Path=../../system/public/type/webfonts/,Extension=.ttf]
\newfontfamily\aclight{ywft-processing-light}[Path=../../system/public/type/webfonts/,Extension=.ttf]
\setmonofont{Latin Modern Mono}[
  Extension=.otf,
  UprightFont=lmmono10-regular,
  ItalicFont=lmmono10-italic,
  Scale=0.85,
]

\usepackage{xcolor}
\usepackage{titlesec}
\usepackage{enumitem}
\usepackage{booktabs}
\usepackage{fancyhdr}
\usepackage{graphicx}
\graphicspath{{figures/}{../../papers/arxiv-ac/figures/}}
\usepackage{ragged2e}
\usepackage{microtype}
\usepackage{eso-pic}
\usepackage{tikz}
\usepackage{hyperref}
\usepackage{xurl}

\definecolor{acpink}{RGB}{180,72,135}
\definecolor{acpurple}{RGB}{120,80,180}
\definecolor{acdark}{RGB}{64,56,74}
\definecolor{acgray}{RGB}{119,119,119}
\definecolor{palspink}{RGB}{205,92,155} % #cd5c9b — canonical purple-pals.svg fill

% === PALS SIDE-STAMPS (mirrors bios.mjs side-stamp pattern + ac-paper-layout.sty) ===
\AddToShipoutPictureBG{%
  \begin{tikzpicture}[remember picture, overlay]
    \pgfmathsetmacro{\palswiggle}{7*sin(\thepage * 67)}%
    \node[opacity=0.10, anchor=center, rotate={90 + \palswiggle}]
      at ([xshift=-1.0cm, yshift=-0.33\paperheight]current page.north west)
      {
\includegraphics[width=13cm]{palsmark}};
    \node[opacity=0.10, anchor=center, rotate={-90 + \palswiggle}]
      at ([xshift=1.0cm, yshift=-0.33\paperheight]current page.north east)
      {
\includegraphics[width=13cm]{palsmark}};
  \end{tikzpicture}%
}

\hypersetup{colorlinks=true,linkcolor=acpurple,urlcolor=acpurple,citecolor=acpurple,
  pdftitle={The Pals Mark: История логотипа Aesthetic.Computer}}

\titleformat{\section}{\normalfont\bfseries\normalsize\uppercase}{\thesection.}{0.5em}{}
\titlespacing{\section}{0pt}{1.2em}{0.3em}
\titleformat{\subsection}{\normalfont\bfseries\small}{\thesubsection}{0.5em}{}
\titlespacing{\subsection}{0pt}{0.8em}{0.2em}

\pagestyle{fancy}\fancyhf{}
\renewcommand{\headrulewidth}{0pt}
\fancyhead[C]{\footnotesize\color{acpink}\textit{The Pals Mark --- \ac{}}}
\fancyfoot[C]{\footnotesize\thepage}

\setlist[itemize]{nosep, leftmargin=1.2em, itemsep=0.15em}
\setlist[enumerate]{nosep, leftmargin=1.4em}
\setlength{\columnsep}{1.8em}
\setlength{\parindent}{1em}
\setlength{\parskip}{0.3em}

% Narrow two-column layout. Russian is set without hyphenation patterns
% here, so long words are unbreakable — give the line-breaker much more
% slack to avoid words spilling into the column gutter.
\tolerance=4000
\emergencystretch=4em
\hyphenpenalty=50
\hbadness=10000

\newcommand{\ac}{\textsc{Aesthetic}{\color{acpink}.}\textsc{Computer}}
\newcommand{\acdot}{{\color{acpink}.}}
\newcommand{\pals}{\textsc{Pals}}
% commit link: \commit{shorthash}{message} -> links the hash to GitHub
\newcommand{\commit}[2]{\href{https://github.com/whistlegraph/aesthetic-computer/commit/#1}{\texttt{#1}}, \emph{``#2''}}
% repo file link
\newcommand{\repofile}[1]{\href{https://github.com/whistlegraph/aesthetic-computer/blob/main/#1}{\texttt{#1}}}

% Cover language bar. The same paper is published in six languages; each
% edition links to the others by canonical permalink. \klanghere marks the
% current language; \klang links a sibling edition by its URL suffix.
\newcommand{\klang}[2]{\href{https://papers.aesthetic.computer/pals-mark-26-arxiv#2.pdf}{\color{acgray}#1}}
\newcommand{\klanghere}[1]{{\bfseries\color{acpink}#1}}
\newcommand{\klangsep}{{\color{acgray!60}\,\textperiodcentered\,}}

\begin{document}

\twocolumn[{%
\begin{center}
{\acbold\fontsize{27pt}{31pt}\selectfont\color{acdark} The \pals{} Mark}\par
\vspace{0.5em}
\includegraphics[width=0.6\textwidth]{figures/cover}\par
\vspace{0.4em}
{\itshape\fontsize{14.5pt}{17pt}\selectfont\color{acpink} История логотипа Aesthetic{\acdot}Computer}\par
\vspace{0.5em}
{\small\color{acgray} \href{https://prompt.ac/@jeffrey}{@jeffrey} \,$\cdot$\, Aesthetic{\acdot}Computer \,$\cdot$\, логотип: Niki Stebbins \,$\cdot$\, ORCID \href{https://orcid.org/0009-0007-4460-4913}{0009-0007-4460-4913}}\par
\vspace{0.45em}
{\footnotesize\sffamily\klang{EN}{}\klangsep\klang{DA}{-da}\klangsep\klang{ES}{-es}\klangsep\klang{ZH}{-zh}\klangsep\klang{JA}{-ja}\klangsep\klanghere{RU}}\par
\vspace{0.55em}
{\color{acpink}\rule{\textwidth}{1.2pt}}
\vspace{0.5em}
\end{center}

\begin{quote}
\small\noindent\textbf{Аннотация.}
\pals{} --- это логотип \ac{}: две нарисованные от руки фигуры, держащиеся за руки с поднятыми вверх руками. Эта статья рассказывает, как был создан знак --- заказанный художнице Niki Stebbins и нарисованный летом 2023 года \emph{внутри} \ac{} с помощью её собственного инструмента whistlegraph \texttt{wgr}, --- и как затем он стал инфраструктурой: вращающимся «граббером логотипа» \texttt{pals.aesthetic.computer}, который выдаёт множество поз знака в аватарах, статьях, нативной ОС и вебе, и каждая из них --- ссылка обратно на проект.
\end{quote}
\vspace{0.5em}
}]

\section{Что такое \ac{}}

\ac{} --- это mobile-first среда исполнения и социальная сеть для
творческих вычислений: единая строка URL, которая одновременно служит
интерфейсом, подобным музыкальному инструменту, где запоминаемые команды
вызывают интерактивные программы, называемые «пьесами». Она поддерживает
как JavaScript (\texttt{.mjs}), так и собственный диалект Lisp ---
\textbf{KidLisp} (118 встроенных функций). Вся система общается с пьесами
через один основной API (\texttt{disk.mjs}, $\sim$600\,КБ).

\textbf{Когда система появилась онлайн.} Домен \texttt{aesthetic.computer}
удерживается с \textbf{2013-12-12}. Платформа в её нынешнем виде выросла из
\emph{No Paint} (2020), разработка началась в \textbf{2021}, и она
непрерывно развёртывается из текущего репозитория с \textbf{декабря 2022
года} ($\sim$20{,}300 коммитов).

\textbf{Статистика платформы} (в реальном времени, \today):
\begin{itemize}
  \item $\sim$400 встроенных пьес в репозитории (383 \texttt{.mjs} + 20
        \texttt{.lisp}); 265 опубликованных пользователями работ.
  \item \textbf{17{,}538} программ KidLisp хранятся в продакшене ---
        растущий корпус генеративного искусства, созданного сообществом.
  \item Федерация ATProto / Bluesky; коллекции записей по каждому хэндлу
        (рисунки, настроения, пьесы, KidLisp) в сети \texttt{at://}.
\end{itemize}

\textbf{Кому она служит и кто знает знак.} Пользователи \ac{} --- это
сообщество творческих кодеров, художников и нетехнических «Pals», которые
мгновенно публикуют работы под пьесами \texttt{@handle}. Платформа
вовлекает широкую аудиторию во всём поле искусства-и-технологии, и работы
\ac{} вошли в институциональные собрания. Сам знак \pals{} едет по каждой
из этих поверхностей: экран загрузки, статьи, аватары в соцсетях, нативная
ОС. Его концептуальная родословная в сообществе радикальной компьютерной
музыки \emph{Goodiepal \& Pals} (см. раздел о родословной ниже) придаёт
имени признанную историю в медиаискусстве за пределами собственной
аудитории \ac{}.

\section{Что изображает знак}

Логотип \pals{} --- это \textbf{две контурные человеческие фигуры,
держащиеся за руки с поднятыми вверх руками}, --- как два друга, идущие
или танцующие вместе, --- выполненные свободной, органичной, нарисованной
от руки линией. Каноническая статическая векторная версия
(\texttt{system/public/purple-pals.svg}) --- розово-пурпурная,
\texttt{\#cd5c9b}. Это \emph{не} «пара глаз» --- раннее ошибочное
прочтение, с тех пор исправленное в наборе статей.

Важно, что знак был создан \emph{с помощью} платформы, а не просто
\emph{для} неё: Niki рисовала эскизы логотипа с помощью \texttt{wgr} ---
пьесы-рисовалки, отрисовщика whistlegraph, работающего внутри самого
\ac{}. Поэтому логотип --- это одновременно и айдентика, и демонстрация
инструмента.

\section{Как был создан знак}

\subsection{Сначала --- штамп}
Всё начинается летом \textbf{2020 года} с небольшого подарка. Niki
Stebbins присылает по почте рисунок --- маленькую фиолетовую звёздочку, ---
и Jeffrey вставляет её в игру как штамп под хэндлом \texttt{@niki} и просит
ещё. Следующие два года сотрудничество представляет собой свободную,
непрерывную \emph{сессию рисования}: штампы, игра-рисовалка, каракули
whistlegraph. К \textbf{2022 году} это уже оплачиваемая дизайнерская
работа. Ничто из этого ещё не логотип, но это почва, из которой логотип
вырастает, --- привычка двух людей рисовать вместе внутри инструмента.

\subsection{Возврат к дизайну (июнь 2023)}
Собственно логотип начинается с неудовлетворённости. В \textbf{июне 2023
года} Niki пишет, что она «не на 100\% довольна логотипом, который я тебе
оставила», и возвращается к нему. Она вслух размышляет о цвете ---
«например, фиолетовый или зелёный\dots{} не синий, слишком много синих
логотипов приложений» --- и о форме: должен ли знак быть «как рамка, только
замкнутая фигура», чем-то, внутри чего можно работать? Она присылает
страницу скетчбука с вариантами и выбирает любимый --- «мне больше всего
нравится обведённый в кружок», --- но тут же беспокоится, что он «довольно
похож на логотип are.na», и продолжает. Ключевая инструкция, которую она
даёт себе сама, --- это исправление пропорций: сделать так, «чтобы он
выглядел меньше как звёзды и больше как фигуры». Страница набросков получает
собственный вердикт --- «похоже на морскую звезду, лол», --- и чтобы это
исправить, она вовсе уходит от экрана, завязывая и фотографируя куски
\textbf{верёвки} в парке, чтобы понять, как на самом деле падают два
связанных тела.

\subsection{Построение двух фигур (июль--август 2023)}
В течение лета знак превращается в \emph{две} фигуры вместо одной. Niki
покупает подержанную цифровую камеру \textbf{Kodak EasyShare} специально
для фотографирования логотипа, задаёт версии «точный размер логотипа VS
Marketplace» и спрашивает, хочет ли Jeffrey «ещё рисунков логотипа в a.c
для аватарок в соцсетях и штуки для app-store». Репозиторий подтверждает
переписку: в \commit{5bcefc5a9}{better defaults in \texttt{wgr} for niki's
logo efforts} сам инструмент рисования настраивается под неё, а к
\commit{a0009f5d0}{add new vscode extension logo from niki} готовый ассет
уже поставлен. В конце августа она сообщает, что «варианты \textbf{двух
логотипов}» готовы, точно описывая жест: «правая фигура тянула
левую\dots{} рука левой фигуры выше». Теперь знак --- это то, чем он
остаётся до сих пор: два человека, руки вверх, один тянет другого за собой.
\textbf{2023-08-31} выбрана картинка профиля в Instagram.

\subsection{Мороженое и запуск (сентябрь 2023)}
Завершение --- самая странная и лучшая часть. Чтобы сфотографировать знак в
реальном мире, Niki рисует его в натуральную величину на асфальте и
пытается снять его с мороженым, --- но мороженое на стадионе продаётся в
стаканчиках, поэтому она заезжает в McDonald's за рожком и делает снимок.
\textbf{2023-09-09} она «отправила логотипы на почту a.c» с темой
\emph{!c3 cr34m l0go\$} (в теле письма: «ice cream logos 8)»). На той же
неделе знак перестаёт быть картинкой и становится инфраструктурой:
появляется \commit{0aba1a378}{add dynamic logo generator}, и Jeffrey
сообщает однострочный итог --- «\^{}\ теперь будет каждый раз возвращать
случайный логотип». Логотип был нарисован \emph{с помощью} платформы (в
инструменте whistlegraph \texttt{wgr}), а затем стал обслуживаться
\emph{ею}. Он никогда не был только картинкой; он всегда был ещё и
программой.

\section{Конечная точка: \texttt{pals.aesthetic.computer}}

Знак стал инфраструктурой в сентябре 2023 года, когда \ac{} развернул
выделенный сервис логотипа и дал ему запоминаемый адрес.

\begin{itemize}
  \item \textbf{2023-09-05} --- \commit{0aba1a378}{add dynamic logo
        generator}. Бэкенд-модуль \repofile{system/backend/logo.mjs}
        помечен меткой \texttt{23.09.05.22.31}.
  \item \textbf{2023-09-12} --- \commit{f1a4f272c}{logo -> pals.aesthetic}
        --- маршрут переименован в айдентику \pals{}.
  \item \textbf{2023-09-13} --- \commit{12e348035}{pals.aesthetic.computer
        url fix} --- конечная точка запущена и доступна на собственном
        поддомене.
\end{itemize}

\noindent Итак, \href{https://pals.aesthetic.computer}{\texttt{pals.aesthetic.computer}}
\textbf{появился онлайн 2023-09-12--13}, через три дня после появления
динамического генератора.

\subsection{Как это работает в качестве «граббера логотипа»}

Конечная точка --- это крошечный вращающийся прокси, а не единственное
фиксированное изображение. Каждый запрос возвращает другую позу \pals{}:

\begin{enumerate}
  \item \href{https://pals.aesthetic.computer}{\texttt{pals.aesthetic.computer}}
        перенаправляет (через \repofile{system/netlify.toml}) на
        \href{https://aesthetic.computer/api/logo}{\texttt{aesthetic.computer/api/logo}}.
  \item \repofile{system/backend/logo.mjs} хранит список \emph{рисунков}
        \pals{} в CDN-бакете \texttt{pals-aesthetic-computer} и возвращает
        случайный при каждом вызове (избегая непосредственного повтора).
  \item Функция согласует содержимое: браузер получает HTML-страницу во
        весь экран на случайном фиолетовом фоне, по нажатию
        переключающуюся на следующую позу (передавая
        \texttt{?previousLogo=}, чтобы она никогда не повторялась); запрос
        \texttt{curl}/\texttt{<img>}/сервера получает сырой \texttt{.png}
        напрямую. Таким образом, один и тот же URL --- это и просмотрщик,
        и ассет для прямых ссылок.
\end{enumerate}

\subsection{Множество поз \pals{}}

Набор поз граббера --- это галерея рисунков \texttt{wgr} середины 2023
года (серия \texttt{painting-2023.7.29}--\texttt{2023.8.21}). Поскольку
логотип --- это \emph{коллекция нарисованных поз}, а не один застывший
глиф, бренд читается как живой: каждая загрузка страницы, обновление
аватара или запрос фавикона может показать две фигуры в чуть иной позе.
Полный набор из двенадцати живых поз воспроизведён в Приложении~A.

\section{Использование на всех площадках}

Поскольку любая площадка может напрямую сослаться на
\href{https://pals.aesthetic.computer}{\texttt{pals.aesthetic.computer}}
(или встроить статический
\repofile{system/public/purple-pals.svg}), единый знак пронизывает всю
экосистему, и каждый его экземпляр --- ссылка обратно на проект:

\begin{itemize}
  \item \textbf{Аватары в соцсетях и Gravatar} --- вращающаяся конечная
        точка предоставляет изображение профиля, которое тихо меняет позу,
        так что одна и та же личность появляется в разных сервисах, всегда
        указывая на \ac{}.
  \item \textbf{Веб и маркетинг} --- фавиконы, шапки документации и
        встроенные теги \texttt{<img>} подтягивают живых \pals{}; нажатие
        переключает позы, превращая сам логотип в небольшой интерактивный
        тизер платформы.
  \item \textbf{Статьи} --- \pals{} едут по полям каждого рабочего
        черновика \ac{} в виде полупрозрачных, покачивающихся от страницы
        к странице боковых штампов (см. эту страницу) и появляются в
        колофонах печатных изданий.
  \item \textbf{Нативная ОС и десктоп} --- \pals{} служат брендингом
        загрузки (заставка Plymouth) в AC Native OS и векторной иконкой в
        трее в AC Electron и MenuBand --- с акцентной подкраской, с
        возможностью бейджа.
  \item \textbf{Печать (в ожидании)} --- в мае 2025 года Jeffrey
        обратился в Casey Rubber Stamps, чтобы вырезать физические штампы
        \pals{} трёх размеров; заказ застопорился и так и не был завершён.
\end{itemize}

\section{Почему «\pals{}»? Линия Goodiepal}

Название --- не украшение; это структурное утверждение, и его вес приходит
извне \ac{}. Как утверждается в статье \ac{}
\emph{Radical Computer Art: Goodiepalian Approaches in Aesthetic
Computer} (\repofile{papers/arxiv-goodiepal/goodiepal.tex}) и в цитируемой
ею статье о платформе, «\pals{}» происходит от \textbf{Goodiepal \&
Pals} --- свободной группы единомышленников вокруг датско-фарерского
музыканта и теоретика Radical Computer Music Goodiepal (Gaeoudji SYGNOK,
р.~1974).

Модель \pals{} --- это политика организации: группа без членства, без
институциональной поддержки и без формальной структуры, скрепляемая общей
практикой и взаимным признанием, а не принадлежностью. Goodiepal преподавал
историю и эстетику электронной музыки в Королевской датской музыкальной
академии и ушёл в 2008 году после конфликта, описанного как наличие «не
того типа друзей», объявив «интеллектуальную войну глупости в современной
компьютерной музыке и медиаискусстве». Jeffrey гастролировал с Goodiepal
\& Pals по Европе в 2018 году, и инструменты, созданные для тех
выступлений, --- прямые предки архитектуры \ac{}.

Так логотип вносит эту политическую историю медиаискусства в бренд. Как
говорится в статье о Goodiepal, назвать логотип «\pals{}» --- и озаглавить
партитуру \emph{Score for Aesthetic{\color{acpink}.}Computer \& Pals} ---
«это не брендинг; это структурное утверждение: платформа принадлежит своему
сообществу практиков, а не разработчику или инвесторам». Две фигуры,
держащиеся за руки, созданные \emph{внутри} инструмента, --- это аргумент,
воплощённый в знаке: \ac{} --- \emph{для} каждого, кто приходит, --- для
\pals{}.

\section{Авторство}

Jeffrey нанял \textbf{Niki Stebbins} для разработки логотипа \ac{} под его
руководством. На протяжении 2022--2023 годов она разрабатывала знак внутри
собственного инструмента рисования платформы \texttt{wgr}, прорабатывая
эскизы, верёвку и мороженое, задокументированные выше, и создала устройство
из двух фигур, поставляемое как \texttt{purple-pals.svg}. Вращающийся
граббер \texttt{pals.aesthetic.computer} был запущен 2023-09-12--13;
каноническая статическая векторная версия впервые появляется 2024-01-31
(\commit{bf29eadf8}{host purple-pals graphic for docs}); а с 2026-02
репозиторий и статьи явно ставят «\pals{}» в центр как айдентику платформы
(\commit{48dd8923e}{rename score to ``Score for
Aesthetic.Computer \& Pals''}). Знак --- это заказная работа: автор ---
Niki, руководитель --- Jeffrey, правообладатель --- \ac{}.

\vspace{0.6em}
{\footnotesize\color{acgray}\noindent Источники: история коммитов
репозитория;
\repofile{system/backend/logo.mjs};
\repofile{system/netlify/functions/logo.js};
\repofile{system/netlify.toml};
\repofile{papers/arxiv-goodiepal/goodiepal.tex};
\repofile{system/public/purple-pals.svg}.\par}

% ============================ APPENDICES ============================
\onecolumn
\appendix

\noindent\emph{Замечание о датах.} Каждый визуальный ассет ниже подписан
датой его создания: позы \texttt{wgr} --- по временной метке,
закодированной в слаге рисунка; ассеты монорепозитория --- по дате их
первого коммита (git \texttt{--diff-filter=A}); живые снимки --- по дате,
когда они были сфотографированы из продакшена.

\section*{Приложение A --- Множество поз (живой набор \texttt{pals.aesthetic.computer})}
\addcontentsline{toc}{section}{Приложение A --- Позы}
\noindent Двенадцать рисунков \texttt{wgr}, которые перебирает граббер
логотипа, подтянутые вживую из CDN \texttt{pals-aesthetic-computer}. Один и
тот же мотив из двух фигур, каждая --- в другой позе и цветовом решении.
\vspace{1.2em}

\begin{minipage}[t]{0.242\textwidth}\centering
  
\includegraphics[width=\linewidth,height=\linewidth,keepaspectratio]{figures/poses/small/pose-1.png}\\[2pt]
  {\scriptsize\textbf{pose 1}}\\[-1pt]
  {\scriptsize\color{acgray}2023-08-01}
\end{minipage}\hfill\begin{minipage}[t]{0.242\textwidth}\centering
  \includegraphics[width=\linewidth,height=\linewidth,keepaspectratio]{figures/poses/small/pose-2.png}\\[2pt]
  {\scriptsize\textbf{pose 2}}\\[-1pt]
  {\scriptsize\color{acgray}2023-08-21}
\end{minipage}\hfill\begin{minipage}[t]{0.242\textwidth}\centering
  
\includegraphics[width=\linewidth,height=\linewidth,keepaspectratio]{figures/poses/small/pose-3.png}\\[2pt]
  {\scriptsize\textbf{pose 3}}\\[-1pt]
  {\scriptsize\color{acgray}2023-08-21}
\end{minipage}\hfill\begin{minipage}[t]{0.242\textwidth}\centering
  \includegraphics[width=\linewidth,height=\linewidth,keepaspectratio]{figures/poses/small/pose-4.png}\\[2pt]
  {\scriptsize\textbf{pose 4}}\\[-1pt]
  {\scriptsize\color{acgray}2023-08-01}
\end{minipage}

\vspace{0.8em}

\begin{minipage}[t]{0.242\textwidth}\centering
  \includegraphics[width=\linewidth,height=\linewidth,keepaspectratio]{figures/poses/small/pose-5.png}\\[2pt]
  {\scriptsize\textbf{pose 5}}\\[-1pt]
  {\scriptsize\color{acgray}2023-08-21}
\end{minipage}\hfill\begin{minipage}[t]{0.242\textwidth}\centering
  \includegraphics[width=\linewidth,height=\linewidth,keepaspectratio]{figures/poses/small/pose-6.png}\\[2pt]
  {\scriptsize\textbf{pose 6}}\\[-1pt]
  {\scriptsize\color{acgray}2023-08-01}
\end{minipage}\hfill\begin{minipage}[t]{0.242\textwidth}\centering
  \includegraphics[width=\linewidth,height=\linewidth,keepaspectratio]{figures/poses/small/pose-7.png}\\[2pt]
  {\scriptsize\textbf{pose 7}}\\[-1pt]
  {\scriptsize\color{acgray}2023-08-21}
\end{minipage}\hfill\begin{minipage}[t]{0.242\textwidth}\centering
  \includegraphics[width=\linewidth,height=\linewidth,keepaspectratio]{figures/poses/small/pose-8.png}\\[2pt]
  {\scriptsize\textbf{pose 8}}\\[-1pt]
  {\scriptsize\color{acgray}2023-08-21}
\end{minipage}

\vspace{0.8em}

\begin{minipage}[t]{0.242\textwidth}\centering
  \includegraphics[width=\linewidth,height=\linewidth,keepaspectratio]{figures/poses/small/pose-9.png}\\[2pt]
  {\scriptsize\textbf{pose 9}}\\[-1pt]
  {\scriptsize\color{acgray}2023-08-01}
\end{minipage}\hfill\begin{minipage}[t]{0.242\textwidth}\centering
  \includegraphics[width=\linewidth,height=\linewidth,keepaspectratio]{figures/poses/small/pose-10.png}\\[2pt]
  {\scriptsize\textbf{pose 10}}\\[-1pt]
  {\scriptsize\color{acgray}2023-08-01}
\end{minipage}\hfill\begin{minipage}[t]{0.242\textwidth}\centering
  \includegraphics[width=\linewidth,height=\linewidth,keepaspectratio]{figures/poses/small/pose-11.png}\\[2pt]
  {\scriptsize\textbf{pose 11}}\\[-1pt]
  {\scriptsize\color{acgray}2023-07-31}
\end{minipage}\hfill\begin{minipage}[t]{0.242\textwidth}\centering
  
\includegraphics[width=\linewidth,height=\linewidth,keepaspectratio]{figures/poses/small/pose-12.png}\\[2pt]
  {\scriptsize\textbf{pose 12}}\\[-1pt]
  {\scriptsize\color{acgray}2023-07-29}
\end{minipage}

\vspace{0.8em}




\newpage
\section*{Приложение B --- Знак по всему монорепозиторию}
\addcontentsline{toc}{section}{Приложение B --- Ассеты монорепозитория}
\noindent Одни и те же \pals{}, нарисованные однажды и переоформленные
повсюду: каноническая векторная версия и её потомки по всей кодовой базе
--- расширение редактора, заставка загрузки нативной ОС, трей десктопа,
обложки pop, штамп на счетах и набор цветовых решений keymaps. Каждая
плитка подписана датой первого коммита.
\vspace{1.2em}

\begin{minipage}[t]{0.194\textwidth}\centering
  
\includegraphics[width=\linewidth,height=\linewidth,keepaspectratio]{figures/assets/c01-purple-pals.png}\\[2pt]
  {\scriptsize\textbf{purple-pals.svg}}\\[-1pt]
  {\scriptsize\color{acgray}2024-01-31}
\end{minipage}\hfill\begin{minipage}[t]{0.194\textwidth}\centering
  
\includegraphics[width=\linewidth,height=\linewidth,keepaspectratio]{figures/assets/c02-vscode.png}\\[2pt]
  {\scriptsize\textbf{vscode-extension}}\\[-1pt]
  {\scriptsize\color{acgray}2024-01-31}
\end{minipage}\hfill\begin{minipage}[t]{0.194\textwidth}\centering
  
\includegraphics[width=\linewidth,height=\linewidth,keepaspectratio]{figures/assets/c03-tvos.png}\\[2pt]
  {\scriptsize\textbf{apple/tvos-tapes}}\\[-1pt]
  {\scriptsize\color{acgray}2024-01-31}
\end{minipage}\hfill\begin{minipage}[t]{0.194\textwidth}\centering
  
\includegraphics[width=\linewidth,height=\linewidth,keepaspectratio]{figures/assets/c04-plymouth.png}\\[2pt]
  {\scriptsize\textbf{AC-OS boot splash}}\\[-1pt]
  {\scriptsize\color{acgray}2026-02-21}
\end{minipage}\hfill\begin{minipage}[t]{0.194\textwidth}\centering
  
\includegraphics[width=\linewidth,height=\linewidth,keepaspectratio]{figures/assets/c05-tray.png}\\[2pt]
  {\scriptsize\textbf{Electron tray}}\\[-1pt]
  {\scriptsize\color{acgray}2026-05-20}
\end{minipage}

\vspace{0.8em}

\begin{minipage}[t]{0.194\textwidth}\centering
  
\includegraphics[width=\linewidth,height=\linewidth,keepaspectratio]{figures/assets/c14-stamp.png}\\[2pt]
  {\scriptsize\textbf{invoice stamp}}\\[-1pt]
  {\scriptsize\color{acgray}2026-03-06}
\end{minipage}\hfill\begin{minipage}[t]{0.194\textwidth}\centering
  
\includegraphics[width=\linewidth,height=\linewidth,keepaspectratio]{figures/assets/c12-gray.png}\\[2pt]
  {\scriptsize\textbf{grayscale (quotes)}}\\[-1pt]
  {\scriptsize\color{acgray}2026-05-21}
\end{minipage}\hfill\begin{minipage}[t]{0.194\textwidth}\centering
  
\includegraphics[width=\linewidth,height=\linewidth,keepaspectratio]{figures/assets/c13-hellsine.png}\\[2pt]
  {\scriptsize\textbf{pop cover glyph}}\\[-1pt]
  {\scriptsize\color{acgray}2026-05-24}
\end{minipage}\hfill\begin{minipage}[t]{0.194\textwidth}\centering
  
\includegraphics[width=\linewidth,height=\linewidth,keepaspectratio]{figures/assets/c06-pink.png}\\[2pt]
  {\scriptsize\textbf{keymaps pink}}\\[-1pt]
  {\scriptsize\color{acgray}2026-06-13}
\end{minipage}\hfill\begin{minipage}[t]{0.194\textwidth}\centering
  
\includegraphics[width=\linewidth,height=\linewidth,keepaspectratio]{figures/assets/c07-blue.png}\\[2pt]
  {\scriptsize\textbf{keymaps blue}}\\[-1pt]
  {\scriptsize\color{acgray}2026-06-13}
\end{minipage}

\vspace{0.8em}

\begin{minipage}[t]{0.194\textwidth}\centering
  
\includegraphics[width=\linewidth,height=\linewidth,keepaspectratio]{figures/assets/c08-green.png}\\[2pt]
  {\scriptsize\textbf{keymaps green}}\\[-1pt]
  {\scriptsize\color{acgray}2026-06-13}
\end{minipage}\hfill\begin{minipage}[t]{0.194\textwidth}\centering
  
\includegraphics[width=\linewidth,height=\linewidth,keepaspectratio]{figures/assets/c09-orange.png}\\[2pt]
  {\scriptsize\textbf{keymaps orange}}\\[-1pt]
  {\scriptsize\color{acgray}2026-06-13}
\end{minipage}\hfill\begin{minipage}[t]{0.194\textwidth}\centering
  
\includegraphics[width=\linewidth,height=\linewidth,keepaspectratio]{figures/assets/c10-purple.png}\\[2pt]
  {\scriptsize\textbf{keymaps purple}}\\[-1pt]
  {\scriptsize\color{acgray}2026-06-13}
\end{minipage}\hfill\begin{minipage}[t]{0.194\textwidth}\centering
  
\includegraphics[width=\linewidth,height=\linewidth,keepaspectratio]{figures/assets/c11-yellow.png}\\[2pt]
  {\scriptsize\textbf{keymaps yellow}}\\[-1pt]
  {\scriptsize\color{acgray}2026-06-13}
\end{minipage}\hfill

\newpage
\section*{Приложение C --- Живые размещения (реальные интерфейсы)}
\addcontentsline{toc}{section}{Приложение C --- Живые размещения}
\noindent \pals{}, сфотографированные живущими в реальных
продакшен-интерфейсах, каждый подписан своим исходным URL (ссылки активны):
страница в Visual Studio Marketplace (\pals{} в роли иконки
продукта-расширения), README монорепозитория на GitHub и на каноническом
хосте Tangled, а также вращающаяся конечная точка
\href{https://pals.aesthetic.computer}{\texttt{pals.aesthetic.computer}}
(тот же URL, разные позы при перезагрузке --- плюс сырой
\texttt{/api/logo.png}, отдаваемый тегам \texttt{<img>} и аватарам).
\vspace{1.0em}

% Hand-authored: the Pals living in real production UIs, each with a caption + live URL.

\newcommand{\uishot}[3]{%
  % #1 image  #2 caption  #3 url
  \begin{center}
  \fbox{\includegraphics[width=0.96\linewidth]{#1}}\\[0.4em]
  {\footnotesize #2}\\[0.1em]
  {\footnotesize\ttfamily\href{#3}{#3}}
  \end{center}
  \vspace{0.7em}
}
\newcommand{\uishottwo}[6]{%
  % (#1 img #2 cap #3 url) (#4 img #5 cap #6 url)
  \noindent
  \begin{minipage}[t]{0.485\linewidth}\centering
    \fbox{\includegraphics[width=0.97\linewidth]{#1}}\\[0.35em]
    {\footnotesize #2}\\[0.1em]{\scriptsize\ttfamily\href{#3}{#3}}
  \end{minipage}\hfill
  \begin{minipage}[t]{0.485\linewidth}\centering
    \fbox{\includegraphics[width=0.97\linewidth]{#4}}\\[0.35em]
    {\footnotesize #5}\\[0.1em]{\scriptsize\ttfamily\href{#6}{#6}}
  \end{minipage}
  \vspace{0.9em}
}

\uishot{figures/live/ui-vscode-marketplace.png}{%
  \textbf{Visual Studio Marketplace} --- the \pals{} as the product icon of
  the official \ac{} VS~Code extension (1{,}201 installs).}{%
  https://marketplace.visualstudio.com/items?itemName=aesthetic-computer.aesthetic-computer-code}

\uishottwo
  {figures/live/ui-github-readme.png}{\textbf{GitHub} --- the \pals{} atop the
   monorepo README.}{https://github.com/whistlegraph/aesthetic-computer}
  {figures/live/ui-tangled-readme.png}{\textbf{Tangled} (canonical git host)
   --- the same README, same mark.}{https://tangled.org/aesthetic.computer/core}

\uishot{figures/compare/pals-live-strip.png}{%
  \textbf{\texttt{pals.aesthetic.computer}} in a live browser --- the same URL
  reloaded shows different poses (the rotating ``logo grabber''); the raw
  \texttt{/api/logo.png} is what \texttt{<img>} tags and avatars (Gravatar,
  social) fetch across the web.}{%
  https://pals.aesthetic.computer}



\end{document}
